\documentclass[12pt, a4paper]{article}
\usepackage[utf8]{inputenc}
\usepackage[spanish]{babel}
\usepackage{amsmath}
\usepackage{amssymb}
\usepackage{geometry}
\geometry{margin=2.5cm}

\title{Análisis completo de cónicas \\ Elipse e Hipérbola}
\author{}
\date{}

\begin{document}

\maketitle

\section{Elipse}
\textbf{Ecuación de la elipse:}
\[
5x^{2} + 8xy + 13y^{2} - 22x - 26y + 11 = 0
\]

Determina todos los elementos posibles de la elipse anterior. Para ello, calcula y encuentra:

\begin{enumerate}
    \item El centro \((h, k)\) de la elipse.
    \item El ángulo de rotación \(\theta\) necesario para eliminar el término \(xy\).
    \item La ecuación canónica (forma estándar) después de la rotación y traslación.
    \item Los semiejes mayor \((a)\) y menor \((b)\).
    \item La longitud del eje mayor y del eje menor.
    \item La excentricidad \(e\).
    \item La distancia focal \(c\) y las coordenadas de los dos focos (en el sistema original y en el sistema rotado).
    \item Los vértices (extremos del eje mayor) y covértices (extremos del eje menor).
    \item La longitud del \textit{latus rectum} (lado recto) y las ecuaciones de los dos \textit{latus recta}.
    \item Las ecuaciones de las dos directrices.
    \item El área de la elipse.
    \item Una aproximación del perímetro (usa al menos dos fórmulas diferentes, como la de Ramanujan).
    \item La ecuación de la tangente en el punto donde \(x = 1\) (si existen).
    \item Verifica si el punto \((2,1)\) está dentro, sobre o fuera de la elipse.
    \item (Bonus difícil): Escribe las ecuaciones paramétricas de la elipse en el sistema rotado.
\end{enumerate}

\vspace{1.5cm}

\section{Hipérbola}
\textbf{Ecuación de la hipérbola:}
\[
5x^{2} + 8xy - 13y^{2} - 22x - 26y + 11 = 0
\]

Determina todos los elementos posibles de la hipérbola anterior. Para ello, calcula y encuentra:

\begin{enumerate}
    \item El centro \((h, k)\) de la hipérbola.
    \item El ángulo de rotación \(\theta\) necesario para eliminar el término \(xy\).
    \item La ecuación canónica (forma estándar) después de la rotación y traslación.
    \item Los semiejes \(a\) (transversal) y \(b\) (conjugado).
    \item La longitud del eje transversal y del eje conjugado.
    \item La excentricidad \(e\).
    \item La distancia focal \(c\) y las coordenadas de los dos focos (en el sistema original y en el rotado).
    \item Los vértices (extremos del eje transversal).
    \item La longitud del \textit{latus rectum} (lado recto) y las ecuaciones de los dos \textit{latus recta}.
    \item Las ecuaciones de las dos directrices.
    \item Las ecuaciones de las dos asíntotas.
    \item Una aproximación del perímetro de un arco (o explica por qué la hipérbola no tiene perímetro finito; usa al menos dos aproximaciones si es posible).
    \item La ecuación de la tangente en el punto donde \(x = 1\) (si existe).
    \item Verifica si el punto específico \((2,1)\) está dentro, sobre o fuera de la hipérbola (analiza el signo de la ecuación).
    \item (Bonus difícil): Escribe las ecuaciones paramétricas de la hipérbola en el sistema rotado.
\end{enumerate}

\end{document}